\subsection{不变微分形式}

令 X 为局部有限维的 \(C^{\infty}\) 类实流形，并令 \((g, x) \mapsto gx\) 为连通紧李群 G 在 X 上的 \(C^{\infty}\) 类左作用律。对于 \(g \in G\)，将 \(\tau_{g}\) 记为 X 的自同构 \(x \mapsto gx\)。将 \(\Omega(X)\) 记为 X 上 \(C^{\infty}\) 类实微分形式的代数（《微分与解析流形，结果》，8.3.1）。

对于任意 \(\xi\) 属于 \(\mathfrak{g}\)，将相应的 X 上向量场记为 \(D_{\xi}\)（第 III 章，第 3 节，第 5 号），并将相应的算子 \(\theta(\xi), i(\xi)\) 记为 \(\Omega(X)\) 上的算子，于是有公式（《微分与解析流形，结果》，8.4.5 和 8.4.7）

\[
\theta (\xi) \omega = d (i (\xi) \omega) + i (\xi) d \omega\tag{15}
\]

\[
\frac {d}{d t} (\tau_ {\mathrm{exp} t \xi} ^ {*} \omega) = \tau_ {\mathrm{exp} t \xi} ^ {*} (\theta (\xi) \omega).\tag{16}
\]

若微分形式 \(\omega \in \Omega(\mathrm{X})\) 满足 \(\tau_g^*\omega = \omega\) 对所有 \(g \in \mathrm{G}\) 成立，则称其为不变的；由公式（16），这等价于 \(\theta(\xi)\omega = 0\) 对所有 \(\xi \in \mathfrak{g}\) 成立。将 X 上不变微分形式的空间记为 \(\Omega(\mathrm{X})^{\mathrm{G}}\)；若 \(\omega \in \Omega(\mathrm{X})^{\mathrm{G}}\)，则有 \(d\omega \in \Omega(\mathrm{X})^{\mathrm{G}}\)，所以 \(\Omega(\mathrm{X})^{\mathrm{G}}\) 是复形 \((\Omega(\mathrm{X}), d)\) 的子复形。

定理 2. 典则嵌入 \(\iota : \Omega(\mathrm{X})^{\mathrm{G}} \to \Omega(\mathrm{X})\) 是复形的同伦等价（《代数》，第 X 章，第 33 页，定义 5）；映射 \(\omega \mapsto \omega^{\sharp} = \int_{\mathrm{G}} \tau_g^* \omega dg\) 是一个同伦等价，在同伦意义下为其逆。特别地，映射 \(\mathrm{H}(\iota) : \mathrm{H}(\Omega(\mathrm{X})^{\mathrm{G}}) \to \mathrm{H}(\Omega(\mathrm{X}))\) 是双射。

由第 4 号的推论 3，映射 \(\omega \mapsto \omega^{\sharp}\) 是从 \(\Omega(\mathrm{X})\) 到 \(\Omega(\mathrm{X})^{\mathrm{G}}\) 的复形态射，并且在子复形 \(\Omega(\mathrm{X})^{\mathrm{G}}\) 上诱导恒等映射；因此，要证明定理，只需构造一个分次同态 \(s: \Omega(\mathrm{X}) \to \Omega(\mathrm{X})\)，其次数为 \(-1\)，使得

\[
\omega^ {\sharp} - \omega = (d \circ s + s \circ d) (\omega) \quad \text { for   all } \omega \in \Omega (\mathrm{X}).\tag{17}
\]

由《积分》第 IX 章第 2 节第 4 号的引理 1 以及第 2 节第 2 号的备注 1，存在正测度 \(d\xi\)，定义在具有紧支集的 \(\mathfrak{g}\) 上，其在指数映射下的像等于 \(dg\)。对于 \(\omega \in \Omega(X)\)，置

\[
s (\omega) = \int_ {0} ^ {1} \left\{\int_ {\mathfrak {g}} \tau_ {\exp t \xi} ^ {*} (i (\xi) \omega). d \xi \right\} d t;
\]

我们需证明公式（17）成立。如同第 4 号推论 1 的证明，我们验证公式

\[
d s (\omega) = \int_ {0} ^ {1} \left\{\int_ {\mathfrak {g}} \tau_ {\exp t \xi} ^ {*} d (i (\xi) \omega). d \xi \right\} d t.
\]

现在由公式（15）和（16）得到等式

\[
\begin{array}{r l} & {d s (\omega) + s (d \omega) = \int_ {0} ^ {1} \left\{\int_ {\mathfrak {g}} \tau_ {\exp t \xi} ^ {*} (d (i (\xi) \omega) + i (\xi) d \omega). d \xi \right\} d t} \\ & {\qquad = \int_ {0} ^ {1} \left\{\int_ {\mathfrak {g}} \tau_ {\exp t \xi} ^ {*} (\theta (\xi) \omega). d \xi \right\} d t} \\ & {\qquad = \int_ {\mathfrak {g}} \left\{\int_ {0} ^ {1} \frac {d}{d t} (\tau_ {\exp t \xi} ^ {*} \omega) d t \right\} d \xi} \\ & {\qquad = \int_ {\mathfrak {g}} (\tau_ {\exp \xi} ^ {*} \omega - \omega) d \xi} \\ & {\qquad = \omega^ {\sharp} - \omega ,} \end{array}
\]

定理 2 得证。

将定理应用于 X = G 且 G 通过左平移作用的情形。回忆（第 III 章，第 3 节，第 13 号，命题 50），将 G 上的微分形式与其在单位元处的值对应，给出从 \(\Omega(\mathrm{G})^{\mathrm{G}}\) 到交错多重线性形式的分次代数 \(\operatorname{Alt}(\mathfrak{g})\) 的同构。利用此同构将 \(\Omega(\mathrm{G})^{\mathrm{G}}\) 与 \(\operatorname{Alt}(\mathfrak{g})\) 识别。于是算子 d 由下式给出（第 III 章，第 3 节，第 14 号，命题 51）

\[
\begin{array}{l} d \omega (a _ {1}, \ldots , a _ {p + 1}) \\ = \sum_ {i <   j} (- 1) ^ {i + j} \omega ([ a _ {i}, a _ {j} ], a _ {1}, \ldots , a _ {i - 1}, a _ {i + 1}, \ldots , a _ {j - 1}, a _ {j + 1}, \ldots , a _ {p + 1}) \end{array}
\]

其中 \(\omega\) 属于 \(\mathrm{Alt}^p (\mathfrak{g})\)，且 \(a_1,\ldots ,a_{p + 1}\) 属于 \(\mathfrak{g}\)。

对于 \(\xi \in \mathfrak{g}\)，令相应的左不变向量场为 \(\mathrm{L}_{\xi}\)（利用 G 通过右平移在自身上的作用定义，参见第 III 章，第 3 节，第 6 号）。算子 \(\theta (\mathrm{L}_{\xi}), i(\mathrm{L}_{\xi})\) 与左平移所定义的 G 在 \(\Omega (\mathrm{G})\) 上的作用交换，因而诱导算子 \(\theta (\xi), i(\xi)\)，它们作用在 \(\Omega (\mathrm{G})^{\mathrm{G}}\) 上；在上述识别下，它们由下列公式表示（《微分与解析流形，结果》，8.3.2 和 8.4.2）

\[
(\theta (\xi) \omega) (a _ {1}, \dots , a _ {p}) = - \sum_ {i} \omega (a _ {1}, \dots , a _ {i - 1}, [ \xi , a _ {i} ], a _ {i + 1}, \dots , a _ {p})
\]

\[
(i (\xi) \omega) (a _ {1}, \ldots , a _ {p - 1}) = \omega (\xi , a _ {1}, \ldots , a _ {p - 1})
\]

其中 \(\omega\) 属于 \(\operatorname{Alt}^p (\mathfrak{g})\)，且 \(a_1,\ldots ,a_p\) 属于 \(\mathfrak{g}\)。

由双不变形式组成的子复形 \({}^{\mathrm{G}}\varOmega(\mathrm{G})^{\mathrm{G}}\)（第 III 章，第 3 节，第 13 号）可与子复形 \(\operatorname{Alt}(\mathfrak{g})^{\mathrm{G}}\) 识别，后者由在伴随表示下不变的 g 上交错多重线性形式组成（即，满足 \(\theta(\xi)\omega=0\) 对所有 \(\xi\in\mathfrak{g}\) 成立）。因此有如下复形交换图

\[
\begin{array}{c c c c c} ^ {\mathrm{G}} \varOmega (\mathrm{G}) ^ {\mathrm{G}} & \longrightarrow & \varOmega (\mathrm{G}) ^ {\mathrm{G}} & \longrightarrow & \varOmega (\mathrm{G}) \\ \Big \downarrow & & \Big \downarrow & & \\ \operatorname{Alt} (\mathfrak {g}) ^ {\mathrm{G}} & \longrightarrow & \operatorname{Alt} (\mathfrak {g}) & & \end{array}\tag{18}
\]

其中水平箭头是典则嵌入，竖直箭头是由映射 \(\omega\mapsto\omega(e)\) 诱导的同构。

推论 1. a) 在图（18）中，所有态射都是同伦等价。b) 令 \(\omega \in \operatorname{Alt}(\mathfrak{g})\)。则 \(\omega\) 属于 \(\operatorname{Alt}(\mathfrak{g})^{\mathrm{G}}\) 当且仅当 \(d\omega = 0\) 且 \(d(i(\xi)\omega) = 0\) 对所有 \(\xi \in g\) 成立。复形 \(\operatorname{Alt}(\mathfrak{g})^{\mathrm{G}}\) 的微分为零。c) 分次向量空间 \(\mathrm{H}(\Omega(\mathrm{G}))\) 同构于 \(\operatorname{Alt}(\mathfrak{g})^{\mathrm{G}}\)。

将定理应用于 G 通过左平移在 G 上的作用（分别应用于 G × G 通过 \(((g,h);x)\mapsto gxh^{-1}\) 在 G 上的作用），得到典则嵌入 \(\Omega(\mathrm{G})^{\mathrm{G}}\to\Omega(\mathrm{G})\)（分别为 \({}^{G}\Omega(\mathrm{G})^{\mathrm{G}}\to\Omega(\mathrm{G})\)）是同伦等价；根据《代数》第 X 章第 34 页的推论，断言 a) 成立。

我们证明 \(b)\)。根据第 III 章第 3 节第 14 号命题 51，有 \(d\alpha = -d\alpha\)，即 \(d\alpha = 0\)，对于微分形式 \(\alpha\)（定义在 \(G\) 上且左、右不变）成立。因此，若 \(\omega \in \mathrm{Alt}(\mathfrak{g})^G\)，则 \(d\omega = 0\)，从而 \(d(i(\xi)\omega) = \theta(\xi)\omega - i(\xi)d\omega = 0\)。反之，若 \(d\omega = 0\) 且 \(d(i(\xi)d\omega) = 0\)，则 \(\theta(\xi)\omega = 0\)。断言 \(c)\) 由 \(a)\) 和 \(b)\) 得出。

注. 考虑子复形 \(\mathrm{Z}(\varOmega(\mathrm{G}))\) 和 \(\mathrm{B}(\varOmega(\mathrm{G}))\)（它们属于 \(\varOmega(\mathrm{G})\)）（《代数》，第 X 章，第 25 页）。由两个形式乘积的微分公式（《微分与解析流形，结果》，8.3.5）可知，\(\mathrm{Z}(\varOmega(\mathrm{G}))\) 是 \(\varOmega(\mathrm{G})\) 的子代数，而 \(\mathrm{B}(\varOmega(\mathrm{G}))\) 是 \(\mathrm{Z}(\varOmega(\mathrm{G}))\) 的理想；因此，外积在 \(\mathrm{H}(\varOmega(\mathrm{G}))\) 上诱导分次代数结构。上述结果给出分次代数的同构 \(\mathrm{H}(\varOmega(\mathrm{G})) \to \mathrm{Alt}(\mathfrak{g})^{\mathrm{G}}\)。

令 H 为 G 的闭子群；将定理 2 应用于 X = G/H。根据第 III 章第 1 节第 8 号命题 17 的推论 1，G/H 上的 G-不变微分形式可与 \(\operatorname{Alt}(\mathrm{T}_{e}(\mathrm{G}/\mathrm{H}))\) 的 H-不变元素识别，即与 \(\operatorname{Alt}(\mathfrak{g})\) 中在 H 下不变且被算子 \(i(\xi)\) 对所有 \(\xi \in \mathrm{L}(\mathrm{H})\) 湮没的元素识别。因此：

推论 2. 设 \(\mathrm{H}\) 为 \(\mathrm{G}\) 的闭子群。

\begin{enumerate}
\def\labelenumi{\alph{enumi})}
\item
  典则嵌入 \(\Omega (\mathrm{G / H})^{\mathrm{G}}\to \Omega (\mathrm{G / H})\) 是同伦等价。
\item
  复形 \(\Omega(\mathrm{G}/\mathrm{H})^{\mathrm{G}}\) 可与 \(\operatorname{Alt}(\mathfrak{g})\) 的如下子复形识别：它由元素 \(\omega\)（属于 \(\operatorname{Alt}(\mathfrak{g})\)）在 H 的伴随表示下不变，且满足 \(i(\xi)\omega = 0\) 对所有 \(\xi \in \operatorname{L}(\operatorname{H})\) 成立的元素组成。如果 H 还是连通的，则该子复形由 \(\omega \in \operatorname{Alt}(\mathfrak{g})\) 组成，且满足 \(\theta(\xi)\omega = 0\) 和 \(i(\xi)\omega = 0\) 对所有 \(\xi \in \operatorname{L}(\operatorname{H})\) 成立。
\end{enumerate}

